%% kapa: exam=1 round=35 year=2024 period=2 form=A subject=accounting subject_ko=회계학 seq=15 number=55
%% source: Q-Net 감정평가사 기출문제 (cst00309) artlSeq=5213229 file=raw/1008/5213229/2교시 문제지 원본.pdf
%% conversion: kordoc 4.24.51 -> tools/gen1.py (auto); review=proofread-sonnet; flags=none
\begin{question}{55}
리스제공자 입장에서 일반적으로 금융리스로 분류될 수 있는 조건이 아닌 것은?
\begin{choices}
  \item 리스기간 종료시점에 기초자산의 소유권을 그 시점의 공정가치에 해당하는 변동 지급액으로 이전하는 경우
  \item 기초자산의 소유권이 이전되지는 않더라도 리스기간이 기초자산의 경제적 내용연수의 상당 부분(major part)을 차지하는 경우
  \item 리스약정일 현재, 리스료의 현재가치가 적어도 기초자산 공정가치의 대부분에 해당하는 경우
  \item 기초자산이 특수하여 해당 리스이용자만이 주요한 변경 없이 사용할 수 있는 경우
  \item 리스이용자가 선택권을 행사할 수 있는 날의 공정가치보다 충분히 낮을 것으로 예상되는 가격으로 기초자산을 매수할 수 있는 선택권을 가지고 있고, 그 선택권을 행사할 것이 리스약정일 현재 상당히 확실한 경우
\end{choices}
\end{question}
\answer{1}
